\section{Synthetic transforms and their tape layout}
\label{sec:synthetic-layouts}

We now implement the normalized Fourier transforms needed for
convolution.  Their entries will be short polynomials in a quotient
where the required roots of unity act by exact signed permutations
of coefficient lists.  We also specify where each frequency entry
is stored, because the opposite transform consumes that same order
after a pointwise product.

The polynomial roots and radix-two butterflies belong to the fast
polynomial-transform method of Nussbaumer and
Quandalle~\cite{NussbaumerQuandalle1978,Nussbaumer1980}; we use the
form developed in~\cite[Section~2.4 and Lemma~3.2]{HarveyHoeven2021}.
The work here is to apply the simultaneous layers on a suitable tape
layout and retain the resulting frequency order through convolution.

For an integer $r\geq 2$ that is a power of two, put
\[
 \mathcal R_r=\mathbb C[y]/(y^r+1),\qquad
 \left\|\sum_{k=0}^{r-1}z_k y^k\right\|_\infty
       =\max_{0\leq k<r}|z_k|.
\]
We always use representatives of degree less than $r$.  Multiplication
by $y^a$, for any integer $a$, permutes the coefficients and changes
some signs; consequently it is an isometry for this norm.  For arrays
over $\mathcal R_r$, the same notation denotes the maximum over all
coefficients in all entries.  A disk array has norm at most one.  A
$p$-bit grid array has real and imaginary coefficient parts in
$2^{-p}\mathbb Z$.  For any array, let $Q_p$ truncate each real and
imaginary coefficient part toward zero to this grid.  Then
\begin{equation}
 \|Q_p z-z\|_\infty<\sqrt2\,2^{-p},\qquad
 \|Q_p z\|_\infty\leq\|z\|_\infty.
 \label{eq:transform-truncate}
\end{equation}

\subsection{Roots, butterflies, and stored frequency order}

For a power of two $t$ dividing $2r$, define $\omega_t=y^{2r/t}$.
For $z\in\mathcal R_r^t$ and $0\leq j<t$, set
\[
 (F_t^-z)_j=\frac1t\sum_{k=0}^{t-1}\omega_t^{-jk}z_k,
 \qquad
 (F_t^+z)_j=\frac1t\sum_{k=0}^{t-1}\omega_t^{jk}z_k.
\]
Every power of $\omega_t$ is an isometry, and each transform entry is
an average of $t$ such images.  Thus both transforms are contractions.
The ring has zero divisors, so we verify character cancellation
directly instead of dividing by $\omega_t^j-1$.

\begin{lemma}[Synthetic character cancellation]
\label{lem:synthetic-characters}
If $j$ is an integer, then
\[
 \sum_{k=0}^{t-1}\omega_t^{jk}
 =\begin{cases}t,&t\mid j,\\0,&t\nmid j.\end{cases}
 \qquad F_t^+F_t^-=t^{-1}I.
\]
\end{lemma}
\begin{proof}
Only the second case needs proof.  Write $j=2^v j_0$ with $j_0$ odd
and $0\leq v<\log_2t$.  Replacing $k$ by $k+t/2^{v+1}$ modulo $t$ permutes
the summands and multiplies their sum by
$y^{rj_0}=-1$.  Hence the sum is zero, since the additive group of
$\mathcal R_r$ has no two-torsion.  Expanding the composite transforms
and applying this identity proves the final formula.
\end{proof}

Write
\[
 H_0=\frac12\begin{pmatrix}1&1\\1&-1\end{pmatrix}.
\]
The length-one procedure is the identity.  For $t\geq2$, the
decimation-in-frequency step pairs $z_k$ with $z_{k+t/2}$ for
$0\leq k<t/2$.  The two outputs of $H_0$ occupy these same slots,
whose high address bit is $b=0$ or $b=1$, respectively.  Multiply
the $b=1$ output by $\omega_t^{-k}$, and then apply the length-$t/2$
procedure separately within the two branches.  This computes
$F_t^-$ because, for $b\in\{0,1\}$,
\begin{equation}
 (F_t^-z)_{2h+b}
 =\frac1{t/2}\sum_{k=0}^{t/2-1}
       (\omega_t^2)^{-hk}\omega_t^{-bk}
       \frac{z_k+(-1)^b z_{k+t/2}}2.
 \label{eq:synthetic-dif}
\end{equation}
To follow the recursion, name the original address-bit slots of an
axis of length $t=2^a$ by $a-1,\ldots,0$, from most to least
significant.  For $0\leq s\leq a$, after $s$ levels the values of
the first $s$ named slots are branch bits $b_0,\ldots,b_{s-1}$.
Equation~\eqref{eq:synthetic-dif} identifies these as frequency bits
$0,\ldots,s-1$, numbered from the least significant bit.  The
remaining slots encode an index in $[0,u)$, where $u=t/2^s$.
When $s<a$, the next level uses slot $h=a-1-s$ to distinguish the
paired inputs and then to hold
$b_s$.  Writing $x_v$ for the bit in a still unprocessed slot $v$,
the index used by its twiddle is
\[
       k=\sum_{v=0}^{h-1}2^v x_v,\qquad 0\leq k<u/2.
\]
Previously written branch bits do not enter $k$.  The same description
holds in each shorter branch by \eqref{eq:synthetic-dif}, which proves
the assertion inductively.

After all $a$ rounds, the named slots from high to low therefore
contain the frequency bits from low to high.  If
$\operatorname{rev}_a(j)$ reverses the $a$ binary digits of $j$, the
entry $(F_t^-z)_j$ occupies slot $\operatorname{rev}_a(j)$.
Let $B_t$ denote this record permutation: it moves the entry at $j$
to $\operatorname{rev}_a(j)$.  The recursion writes its outputs
directly in this order; it does not move whole frequency indices.

The procedure consuming this stored order traverses the recursion in
reverse: first process the two smaller branches, then apply the inverse
monomial to the second branch, and finally apply $H_0$.  Each forward
level is invertible, since both its butterfly and its monomial are
invertible.  As $H_0^2=I/2$, the reversed level is half the inverse
of the corresponding forward level.  There are $a=\log_2t$ levels,
so the reverse procedure composed with $B_tF_t^-$ is
$2^{-a}I=I/t$.  Lemma~\ref{lem:synthetic-characters} gives the same
composite for $F_t^+B_t^{-1}$ with $B_tF_t^-$.  The forward procedure
is invertible, so these two reverse operators agree.  Thus the two
exact procedures implement
\begin{equation}
       B_tF_t^-,\qquad F_t^+B_t^{-1},
 \label{eq:stored-opposite}
\end{equation}
respectively.

For $D$ axes of lengths $t_1,\ldots,t_D$, put
\[
 M=\prod_{i=1}^D t_i,\qquad
 F_{\boldsymbol t}^\pm=\bigotimes_{i=1}^D F_{t_i}^\pm,
 \qquad B=\bigotimes_{i=1}^D B_{t_i}.
\]
A round applies one current $H_0$ on each participating axis.  The
twiddle for axis $i$ depends only on that axis's address bits.
The two entries paired by another axis's butterfly have the same
axis-$i$ address, so that twiddle commutes with the other butterfly.
Consequently the individual axis steps may be grouped as a parallel
butterfly $\mathcal H$, followed by one recordwise monomial operation
$J_E$.  At each address, $J_E$ multiplies the record by
\begin{equation}
       y^{E},\qquad
 E=-\sum_{i\ \mathrm{participating}}
       \frac{2r}{u_i}\,b_i k_i\pmod{2r}.
 \label{eq:combined-twiddle}
\end{equation}
Here $u_i$ is that axis's current recursive length, $b_i$ its newly
written branch bit, and $0\leq k_i<u_i/2$ is encoded by its still
unprocessed named slots.  The exponent depends on the address, not
on coefficient values.  In the opposite direction the grouped round
is $\mathcal HJ_E^{-1}$: apply $y^{-E}$ first and the parallel
butterfly second.  Each $\mathcal H$ is a contraction and each
$J_E$ is an isometry, so every completed round is a contraction.

\subsection{A layout on which the parallel layers apply}

To apply Proposition~\ref{prop:simultaneous-layer} in a transform round,
we must place the selected bits of the participating axes at a common
offset in consecutive $K$-bit chunks. We now construct that layout,
keeping each polynomial's coefficients in a contiguous suffix. The
layout uses the single-bit moves of
Lemma~\ref{lem:elementary-streams} and the chunk interchanges of
Lemma~\ref{lem:chunk-swap}; each movement is included in the transform
cost. The layer then supplies one truncated contraction per round,
and its temporary rows and work streams are removed before the next
round begins.

\begin{lemma}[Transform layout and cost]
\label{lem:synthetic-transform-cost}
Fix the rational parameters and positive comparison constants
$a_d\leq b_d$, $a_r\leq b_r$ of
Proposition~\ref{prop:simultaneous-layer}. Fix constants $C_\ell>0$
and $C_w>1$, and choose the layer constants once as
\[
 C_M=C_\ell+1,\qquad
 C_{\rm desc}=32(C_\ell+C_w+1).
\]
Use these constants and $C_w$ in
Proposition~\ref{prop:simultaneous-layer}. These fixed choices determine
a cutoff $p_1$, one fixed-tape procedure, and an implicit time constant
with the following guarantee. For every choice of positive integers
$p,d,r$ with $p\geq p_1$ and
\[
 a_dp^\epsilon\leq d\leq b_dp^\epsilon,\qquad
 2^{a_rp/d}\leq r\leq2^{b_rp/d},
\]
suppose $d\geq2$ and $r=2^\ell$ with an integer $2\leq\ell<p$, and require
\[
 d\ell\leq C_\ell p,\qquad
 1\leq K=\lfloor d^c\rfloor\leq\ell-1.
\]
Let
$t_i\in\{2^{\ell-1},2^\ell\}$ for $1\leq i\leq d-1$, put
$M=\prod_{i=1}^{d-1}t_i$ and $T=rM$, and let $w$ be an integer with
$p+2\leq w\leq C_wp$. The cutoff and time constant are independent of
$p,d,r$, these axis choices, the round headers constructed from them,
and $w$ within the displayed bounds.

The input supplies the self-delimiting header
\[
 \mathcal T=\Gamma(p)\Gamma(d)\Gamma(r)\Gamma(K)\Gamma(w)
       \prod_{i=1}^{d-1}\Gamma(t_i),
\]
where the product is concatenation in axis order and $\Gamma$ is the
integer code of \eqref{eq:integer-descriptor-code}.
The supplied $d$ determines the number of axis entries, and
$\ell=\log_2r$ is obtained from the supplied power of two $r$.
Assume that the header is valid with the stated values and that the
input records use the two's-complement component format of
Proposition~\ref{prop:simultaneous-layer} with width $w$.

The procedure constructs an explicit positional
layout permutation $L$ of record addresses, preserving their polynomial
suffixes, and transforms disk grid arrays of $M$ contiguous
degree-less-than-$r$ polynomial records. Its exact counterpart is
$LB F_{\boldsymbol t}^-$.  Its opposite procedure consumes
that order and has exact counterpart
$F_{\boldsymbol t}^+B^{-1}L^{-1}$.  Each procedure has error less than
$\sqrt2\ell\,2^{-p}$, returns disk grid coefficients in the same
component format, and costs
\begin{equation}
 O\left(Tp\left[dK+pK^{\tau-1}+\ell d^{\lambda'}\right]\right).
 \label{eq:synthetic-transform-cost}
\end{equation}
Both layout construction and final restoration are charged in this
bound.  The forward procedure may retain its output layout for a
pointwise product and the following opposite procedure.
\end{lemma}
\begin{proof}
Put $D=d-1$ and write $a_i=\log_2t_i$. The supplied encoding has
$2w=\Theta(p)$ bits per coefficient, so the logical bit volume is
$\Theta(Tp)$.  Call an axis long if $a_i=\ell$ and short if
$a_i=\ell-1$.  Name the original bit slots of axis $i$ by
$(i,h)$, $a_i-1\geq h\geq0$.  Initially these slots are in axis order
and, within each axis, in decreasing $h$.  The coefficients of each
polynomial form a contiguous suffix.  Put
\[
       b=(\ell-1)\bmod K,
       \qquad q=\lfloor(\ell-1)/K\rfloor.
\]
Form a list of the slots to move to a prefix.  First list
$(i,\ell-1)$ for the length-$2^\ell$ axes in increasing $i$.
Then list $(i,h)$ for $h=b-1,\ldots,0$, and within each $h$ in
increasing $i$; this second list is empty if $b=0$.  Move each listed
slot immediately after the prefix already formed, leaving all other
slots in their relative order.  There are at most $D(b+1)=O(dK)$
listed slots, so the moves cost $O(TpdK)$.  A descriptor records the
physical position of every named slot.  The part of each axis left
outside the prefix now consists of $q$ consecutive $K$-bit chunks.

Number the axes by $i$ and their chunks, from high to low, by $j$.
The remaining physical order changes from
\[
 (1,1),(1,2),\ldots,(1,q),(2,1),\ldots,(D,q)
 \quad\hbox{to}\quad
 (1,1),(2,1),\ldots,(D,1),(1,2),\ldots,(D,q).
\]
To obtain this order, scan its target positions from left to right.
At each position except the last, exchange the chunk currently there
with the required chunk if they differ.  This uses at most $Dq-1$
exchanges of width-$K$ address chunks, even when the chunks are
nonadjacent.  Each costs $O(TpK^\tau)$, so all these exchanges cost
\[
 O\bigl(Tp\,DqK^\tau\bigr)
       =O\bigl(Tp\,pK^{\tau-1}\bigr).
\]
These prefix moves and chunk exchanges define $L$.  We use the same
letter for its induced record permutation: a record at address $x$
moves to address $Lx$, while its polynomial suffix stays intact.
Recording the inverse operations in reverse order gives $L^{-1}$
with the same cost.

For a small illustration of this positional permutation, take
$\ell=6$, $K=2$, and one long axis followed by one short axis.
Then $b=1$ and $q=2$. Write $[i:h,h']$ for the chunk containing
the two named slots $(i,h),(i,h')$, in that order. After $L$ the slots are
\[
 \underbrace{(1,5),(1,0),(2,0)}_{\text{prefix}} ;\qquad
 [1:4,3]\ [2:4,3]\ [1:2,1]\ [2:2,1].
\]
For example, the round at named position $h=4$ uses offset one in
each of the first two chunks. It writes frequency bit one on the long
axis and frequency bit zero on the short axis. The named position
specifies the same physical pairing rule on both axes even though the
frequency significance differs. This example illustrates only the slot
permutation; the parameter bounds of the lemma govern its asymptotic use.

The header has $O(p)$ bits: its axis entries use
$O(\sum_i a_i+d)=O(p)$ bits, and its remaining entries have
$O(p)$ bits in total. The layout lists have binary length polynomial
in $p$. Parsing the header and generating these lists by elementary
integer arithmetic take
$\operatorname{poly}(p)=O(Tp)$ time. Indeed the fixed comparison bands
give $T\geq r\geq2^{(a_r/b_d)p^{1-\epsilon}}$, uniformly over the
permitted $d,r$, so $T$ dominates every fixed power of $p$.

First process the round at named position $h=\ell-1$ on the long
axes, if there are any.  Every axis then has current recursive
length $2^{\ell-1}$.  Next process $h=\ell-2,\ldots,b$ in the $qK$
rounds on the remaining chunks, and finally the bottom positions
$h=b-1,\ldots,0$, if $b>0$.  Thus all participating axes in a
common round have current length $u_i=2^{h+1}$.
The branch bit written at position $h$ is frequency bit $a_i-1-h$;
its frequency significance can differ between long and short axes.
The descriptor locates that named slot after $L$, so the physical
move has not changed which pair the current butterfly uses.  The
opposite procedure executes this schedule in reverse order and,
when present, ends with the extra leading round on the long axes.

For a position $h$ in the remaining chunks, the selected bit has
offset $(h-b)\bmod K$, counted from the least significant bit, in
each of the $D$ consecutive chunks with the same chunk index $j$.
Let $n_P$ and $n_S$ be the numbers of record-address bits before and
after that group.  Every combination of those bits occurs, so the
layer sees the complete address set
\[
 [P]\times[2^K]^D\times[S],\qquad
 P=2^{n_P},\quad S=2^{n_S},\quad
 n_P+DK+n_S=\sum_i a_i\leq D\ell=O(p).
\]
Thus its record count is $PS2^{DK}=M$ and
$\log_2(2M)\leq(C_\ell+1)p=C_Mp$.
The current transform parameters supply $p,d,r,K,w$, and the round
determines $D,P,S$ and the selected offset $\rho$. Encode these values
as the header $\mathcal E$ of Proposition~\ref{prop:simultaneous-layer}.
The same fixed descriptor bound covers every axis choice and round.
Indeed $|\Gamma(v)|\leq2\log_2v+1$, and
$\ell+n_P+n_S\leq d\ell\leq C_\ell p$. Also
$d,D\leq\max\{C_\ell,1\}p$, $K,\rho+1\leq p$, and $w\leq C_wp$.
Adding the nine code lengths therefore gives
\[
 |\mathcal E|
 \leq 2C_\ell p+12\log_2p+
        4\log_2\max\{C_\ell,1\}+2\log_2C_w+9
 \leq C_{\rm desc}p.
\]
The last inequality uses $p\geq1$, $\log_2p\leq p$, and
$\log_2x\leq x$ for $x\geq1$.
This is the descriptor passed to the layer; the longer
table locating named slots remains separate for the transform's address
calculations.
The simultaneous-layer procedure therefore applies at cost
$O(Tpd^{\lambda'})$ per round and $O(Tp\ell d^{\lambda'})$ in total.

In a round on prefix slots, process the participating bits individually.
Temporarily move one such bit immediately before the polynomial
suffix.  In each now adjacent pair, buffer the first record and scan
it with the second, writing the half-sums and half-differences to two
work tapes.  Append the half-sum record and then the half-difference
record, and reverse the bit move after
all pairs are processed.  Including head returns, the work is
$O(rp)$ per pair and $O(Tp)$ for the moves and scans at coefficient
width $O(p)$.
That width suffices without truncating between axes: after $s\leq D$
individual $H_0$ factors, the coefficients are exact dyadics on the
$2^{-(p+s)}$ grid and still lie in the disk.  A temporary sum has
modulus at most two and needs only a constant number of extra integer
bits.  Hence
$p+D+O(1)=O(p)$ bits per component suffice, since $D\leq d$ and
$d\ell=O(p)$.  Complete the combined monomial exactly and apply
$Q_p$ once at the end of the round. Write the resulting disk-grid
numerators in the supplied $w$-bit component format. Their magnitudes
are at most $2^p$, so this takes one linear scan within the same bound.
In an opposite round on prefix slots the
inverse monomial comes first and preserves the grid and disk, so the
same width bound holds for its individual factors.  There are at
most $b+1\leq K$ rounds on prefix slots, giving cost $O(TpdK)$.

For each output record of a forward butterfly round, the descriptor
locates the new branch bit $b_i$ and the lower named slots that encode
$k_i$.  Read their current values to compute
\eqref{eq:combined-twiddle}.  Before an opposite round at position
$h$, all lower positions have already been processed in reverse.
Those slots again index the lower inputs of that level, while slot
$h$ still gives $b_i$.  They therefore encode the same $k_i$ needed
by the inverse monomial before its
butterfly.  The binary descriptions of the addresses, exponent, and
layout have polynomial length in $p$.  Ordinary
fixed-tape arithmetic therefore computes the exponent in
$\operatorname{poly}(p)$ time per polynomial record.  Each record
contains $\Theta(rp)$ bits, and $r$ dominates every fixed power of
$p$, so this work fits within its scan cost.  The same bound covers
scans of the polynomially many round and layout descriptors.
Growing lists reside on tapes; a fixed number of counter and work
tapes suffices for every number of axes.

To implement the monomial, write $E=er+s$ with $0\leq s<r$.
Split the coefficient list before position $r-s$, put its tail
before its head, negate the moved tail, and apply the common sign
$(-1)^e$.  This is multiplication by $y^E$ modulo $y^r+1$.
There are always two pieces.  Splitting them to two tapes and merging
in that order costs $O(rp)$ per record, including head returns;
$s=0$ is an ordinary signed copy.  The inverse twiddle uses the
same split and merge with exponent $-E$, and no coefficient sort is
required.

For a forward round on a group of $D$ chunks, the layer call returns
$Q_p(\mathcal H z)$ and the next operation is $J_E$.  Truncation
toward zero is odd and acts separately on every real and imaginary
part.  Since $J_E$ is a signed coefficient permutation, it follows
that
\[
       J_EQ_p(\mathcal H z)=Q_p(J_E\mathcal H z).
\]
Thus the actual forward call sequence returns exactly the result of
truncating once after the exact contraction round.  In the opposite
direction, $J_E^{-1}$
first preserves the disk grid, and the layer call returns
$Q_p(\mathcal HJ_E^{-1}z)$ directly.  The rounds on prefix slots
have the same forms, with one truncation, by their exact evaluation
above.  A contraction does not increase the norm of the difference
between the computed and exact arrays; the single truncation then
adds less than $\sqrt2\,2^{-p}$ by \eqref{eq:transform-truncate}.
Both the contraction and truncation preserve the disk.  Starting with
zero error, at most $\ell$ rounds therefore
give the claimed error and disk bounds.  The layout permutations
are isometries and add no numerical error.

It remains to identify the order of the full procedures.  Before
relocating slots, write
$\mathcal D_-=BF_{\boldsymbol t}^-$ and
$\mathcal D_+=F_{\boldsymbol t}^+B^{-1}$ for the exact logical forward and
opposite procedures.  Executing each operation on named slots after $L$
conjugates it by $L$.  The forward procedure first moves the records
by $L$, whereas the opposite procedure restores them by $L^{-1}$
after its reverse arithmetic.  Their full exact operators are
\[
 (L\mathcal D_-L^{-1})L=LB F_{\boldsymbol t}^-,
 \qquad
 L^{-1}(L\mathcal D_+L^{-1})
       =F_{\boldsymbol t}^+B^{-1}L^{-1}.
\]
In particular, the forward entry with frequency address $j$ is
stored at $L(Bj)$.  The reverse arithmetic consumes the frequency
branch coordinates and outputs the normalized opposite transform
indexed by the original axis coordinates, which are still in layout $L$.
The factor $B^{-1}$ is realized by that arithmetic; the final
charged movement $L^{-1}$ restores axis-major physical order.
\end{proof}

Pointwise multiplication commutes with applying the same record
permutation $LB$ to both operands.  We may therefore multiply the
forward outputs in their retained order and pass that order directly
to the opposite procedure.  The following polynomial product
implements this operation with an established integer multiplier.

\section{Exact normalized polynomial products}
\label{sec:exact-ring-products}

To multiply the polynomial records, we use the fixed-tape integer
multiplier of Harvey and van der
Hoeven~\cite[Theorem~1.1]{HarveyHoeven2021}, with cost
$O(N\log(2N))$ on $N$-bit inputs.  This is an independent subroutine.
Its inputs below have $O(rp)$ bits, and we include the conversions
between coefficient streams and signed integers in the cost. The
reduction is a signed form of Kronecker substitution; compare
\cite[Lemma~2.5]{HarveyHoeven2021}. We specify the packing and recovery
of coefficients to include negative values and the required tape format.

\begin{lemma}[Signed packing and normalized ring multiplication]
\label{lem:signed-ring-product}
Let $r=2^\ell<2^p$ with integers $\ell\geq1$ and $p\geq2$, and fix $C_w>1$.
Let $w$ be an integer with $p+2\leq w\leq C_wp$.
The input supplies $\Gamma(p)\Gamma(r)\Gamma(w)$ and two disk grid
polynomials $f,g\in\mathcal R_r$ in the component format of
Proposition~\ref{prop:simultaneous-layer} with width $w$.
One can compute the exact Gaussian integer
numerators of $fg$, and then return $Q_p(fg/r)$ in that same format, in time
$O(rp\log(rp))$ on a fixed number of tapes.  The output is a disk
grid polynomial with error less than $\sqrt2\,2^{-p}$ from $fg/r$.
\end{lemma}
\begin{proof}
The supplied power of two $r$ determines $\ell=\log_2r$.
Write
\[
 f=2^{-p}(a+ib),\qquad g=2^{-p}(c+ie),
\]
where $a,b,c,e\in\mathbb Z[y]$ have degree less than $r$ and
coefficients of absolute value at most $2^p$.  Set $B_0=2^{4p}$
and, for $u\in\{a,b,c,e\}$, define the signed evaluation integer
\[
       A_u=u(B_0)=\sum_{j=0}^{r-1}u_jB_0^j.
\]
The bound
\[
 |A_u|\leq2^p\frac{B_0^r-1}{B_0-1}<B_0^r/2
\]
shows that each magnitude fits in $4pr$ bits.  Multiply the four
pairs of absolute values using the established multiplier and restore
each product's sign.  The resulting integers are
$A_aA_c,A_bA_e,A_aA_e,A_bA_c$.  For example, if
$h=ac$ is the ordinary product in $\mathbb Z[y]$, then
\[
       A_aA_c=(ac)(B_0)=\sum_{j=0}^{2r-2}h_jB_0^j.
\]
The other three evaluation products have the same form.  Every
coefficient in any of these ordinary polynomial products satisfies
\begin{equation}
      |h_j|\leq r2^{2p}<2^{3p}<B_0/2.
 \label{eq:packing-separation}
\end{equation}
After the lower coefficients have been removed from an evaluation,
this strict bound identifies the next coefficient by its centered
residue modulo $B_0$.

To form $A_u$, start with carry $\gamma_0=0$ and write
\[
 d_j\equiv u_j+\gamma_j\pmod{B_0},\quad 0\leq d_j<B_0,
 \qquad \gamma_{j+1}=\frac{u_j+\gamma_j-d_j}{B_0}.
\]
Since $-B_0<u_j+\gamma_j<B_0$ whenever
$\gamma_j\in\{-1,0\}$, the carry stays in $\{-1,0\}$.  Telescoping
the recurrence gives
\[
       A_u=\sum_{j=0}^{r-1}d_jB_0^j+\gamma_rB_0^r.
\]
The $r$ blocks are therefore a fixed-width two's-complement encoding
of $A_u$, with $\gamma_r=-1$ exactly when $A_u<0$.  In that case
complement and increment the blocks to obtain $|A_u|$; otherwise
retain them.  Every block has $4p$ bits and every carry has constant
size, so a constant number of scans costs $O(rp)$.  Endianness
changes required by the multiplier are linear-time reversals on a
work tape.

Now let $Z$ be one of the four signed product integers, and let $h$
be its corresponding ordinary polynomial product.  Put $J=2r$ and
pad that product by $h_{J-1}=0$.  The bound on the evaluation integers
gives $|Z|<B_0^J/4$, so $Z$ fits in $J$ blocks of width $4p$.
Convert its sign and magnitude to sign-extended two's complement:
\[
 Z=\sum_{j=0}^{J-1}q_jB_0^j+\varepsilon B_0^J,
 \qquad 0\leq q_j<B_0,\quad \varepsilon\in\{-1,0\}.
\]
Zero-fill the unused high blocks; for a negative sign, complement
and increment within this fixed width.  These conversions take a
constant number of scans.
Put $\rho_0=0$.  At block $j$, let $\widehat h_j$ be the unique
integer in $[-B_0/2,B_0/2)$ congruent to $q_j+\rho_j$ modulo $B_0$,
and put
\[
       \rho_{j+1}=
          \frac{q_j+\rho_j-\widehat h_j}{B_0}.
\]
The carry belongs to $\{0,1\}$: if $\rho_j$ has this property, then
$0\leq q_j+\rho_j\leq B_0$, and centering either leaves this value
unchanged or subtracts $B_0$.

The meaning of the carry is the following residual identity, after
coefficients below $j$ have been correctly extracted:
\[
 \sum_{k=j}^{J-1}h_kB_0^{k-j}
 =\sum_{k=j}^{J-1}q_kB_0^{k-j}
       +\rho_j+\varepsilon B_0^{J-j}.
\]
At $j=0$ this is the equality between the two representations of
$Z=h(B_0)$, with
$\rho_0=0$.  For $j<J$, reduction modulo $B_0$ gives
$h_j\equiv q_j+\rho_j\pmod{B_0}$.  By
\eqref{eq:packing-separation}, including the padded value
$h_{J-1}=0$, the unique centered representative is $h_j$ itself.
Thus $\widehat h_j=h_j$.  Subtracting this coefficient and dividing
the residual identity by $B_0$ gives the next identity, with exactly
the displayed update for $\rho_{j+1}$.  At $j=J$ the identity is
$0=\rho_J+\varepsilon$.  Hence a negative product has final carry
$1$, canceled by its sign extension $\varepsilon=-1$; a nonnegative
product has final carry zero.  The machine implements the induction
by reading one block and updating this constant-size carry at each
step, so extraction costs $O(rp)$.

The real and imaginary numerator streams are respectively
$ac-be$ and $ae+bc$.  To reduce them modulo $y^r+1$, subtract the
coefficient at $j+r$ from the one at $j$ for $0\leq j<r$, treating
a missing coefficient as zero.  Splitting the two halves to tapes
aligns these terms in one further scan.  For each first input index
$j_1\in\{0,\ldots,r-1\}$, exactly one second index $j_2$ in that
range satisfies $j_1+j_2\equiv j\pmod r$.  Its sign is positive
when $j_1+j_2=j$ and negative when $j_1+j_2=j+r$.  Thus each
resulting Gaussian integer coefficient $H_j$ is a sum of exactly
$r$ signed products of Gaussian integer input coefficients of
modulus at most $2^p$.  Consequently
\[
       |H_j|\leq r2^{2p}.
\]
The same triangle inequality before imposing the disk bound gives
$\|fg/r\|_\infty\leq\|f\|_\infty\|g\|_\infty$; in particular the
normalized product of disk polynomials lies in the disk.  All
intermediate coefficient additions fit in $O(p)$ bits.  The
width-$4p$ packing fields and the exact multiplier output prevent
overflow or ambiguity between adjacent fields.

Finally the $p$-bit grid numerator of the normalized product is
\[
 \operatorname{trunc}_0\left(\frac{\Re H_j}{2^{p+\ell}}\right)
   +i\operatorname{trunc}_0\left(\frac{\Im H_j}{2^{p+\ell}}\right).
\]
To implement it, shift the absolute magnitude right by $p+\ell$
bits and restore its sign, at cost $O(p)$ per coefficient.  The returned
coefficient is $2^{-p}$ times this numerator, namely
$Q_p(2^{-2p}H_j/r)$; \eqref{eq:transform-truncate} proves its error
and disk bounds. Each returned component numerator has magnitude at
most $2^p$, so write it as a signed $w$-bit two's-complement word,
most significant bit first, with the real word before the imaginary
word and coefficient indices in increasing order. This costs
$O(rw)=O(rp)$ and supplies the format used by the opposite transform.
Four exact integer products dominate the linear
packing, extraction, combination, wrapping, and truncation costs.
\end{proof}

\subsection{The normalized convolution interface}

For arrays on $\prod_i\mathbb Z/t_i\mathbb Z$ with values in
$\mathcal R_r$, let $*$ denote cyclic convolution in the displayed
axes and define
\[
       \mathcal M_r(f,g)=\frac{f*g}{rM}=\frac{f*g}{T}.
\]
For a fixed output coefficient, each of the $M$ first cyclic
addresses determines one second address, and each of the $r$ first
polynomial indices determines one second index modulo $r$.  The
negacyclic wrap supplies the sign.  Thus that coefficient averages
$rM$ signed scalar products, and $\mathcal M_r$ is a bilinear
contraction:
\[
 \|\mathcal M_r(f,g)\|_\infty
       \leq\|f\|_\infty\|g\|_\infty.
\]
If the other operand is a disk array, replacing one operand changes
the output by at most the norm distance between its two values.

\begin{proposition}[Convolution from the retained transform layout]
\label{prop:synthetic-convolution}
Under the hypotheses of Lemma~\ref{lem:synthetic-transform-cost},
$\mathcal M_r$ on disk grid operands has an approximation with error
less than
\begin{equation}
       \sqrt2\,M(3\ell+1)2^{-p}
 \label{eq:synthetic-convolution-error}
\end{equation}
and cost
\[
O\left(Tp\left[dK+pK^{\tau-1}+\ell d^{\lambda'}
                    +\log(rp)\right]\right).
\]
Put $m_M=\log_2M$ and $w_M=w+m_M$. The returned grid array is in the
original array and coefficient order, with the same two's-complement
component format at width $w_M$. The procedure supplies
$\Gamma(w_M)$ on the output descriptor tape, so the caller knows every
component boundary from this width and the retained shape parameters.
If one operand has norm at most $\rho<1$ and the bound in
\eqref{eq:synthetic-convolution-error} is at most $1-\rho$, the
returned grid array is a disk array.
\end{proposition}
\begin{proof}
Perform both forward transforms and retain their common $LB$ order.
Pass the supplied $p,r,w$ to
Lemma~\ref{lem:signed-ring-product} for corresponding polynomial records.
Its output keeps that record order and uses the component format
required by the opposite procedure. Apply the opposite procedure
to that order, and multiply its result exactly by $M=T/r$.

To check normalization, expansion of the finite sums gives
\[
 F_{\boldsymbol t}^-(f*g)
       =M(F_{\boldsymbol t}^-f)(F_{\boldsymbol t}^-g).
\]
As $F_{\boldsymbol t}^+F_{\boldsymbol t}^-=I/M$, it follows that
\[
 M F_{\boldsymbol t}^+
       \left(\frac{(F_{\boldsymbol t}^-f)
                          (F_{\boldsymbol t}^-g)}r\right)
       =\frac{f*g}{rM}.
\]
Let $P(u,v)$ denote the exact pointwise product $uv/r$, and put
$\mathcal T_+=F_{\boldsymbol t}^+B^{-1}L^{-1}$.
Since $LB$ permutes records and leaves their polynomial suffixes
intact, $P(LBu,LBv)=LBP(u,v)$.  Thus $\mathcal T_+$ applied to the
pointwise product of the exact stored forward outputs equals
$F_{\boldsymbol t}^+P(F_{\boldsymbol t}^-f,F_{\boldsymbol t}^-g)$.
The displayed normalization identity identifies its exact scaling
by $M$ with $\mathcal M_r(f,g)$.

For the error estimate, put
$\delta=\sqrt2\ell\,2^{-p}$ and $\eta=\sqrt2\,2^{-p}$.
Let $u,v$ be the exact stored forward outputs and
$\widetilde u,\widetilde v$ the returned ones.  All four are disk
arrays, and each forward error is less than $\delta$.
The exact map $P$ is bilinear and satisfies
$\|P(x,y)\|_\infty\leq\|x\|_\infty\|y\|_\infty$ by the coefficient
bound in Lemma~\ref{lem:signed-ring-product}.  Hence
\[
\begin{split}
 P(\widetilde u,\widetilde v)-P(u,v)
   &=P(\widetilde u-u,\widetilde v)+P(u,\widetilde v-v),\\
 \|P(\widetilde u,\widetilde v)-P(u,v)\|_\infty
   &<2\delta,
\end{split}
\]
where the first term uses the disk bound on $\widetilde v$ and the
second uses it on $u$.  Write
$\widetilde z=Q_p(P(\widetilde u,\widetilde v))$ and $z_*=P(u,v)$.
The product lemma makes $\widetilde z$ a disk grid array and gives
$\|\widetilde z-z_*\|_\infty<2\delta+\eta$.
If $\widetilde h$ is the returned opposite output, its approximation
bound applies to $\widetilde z$.  Since $\mathcal T_+$ is a contraction,
\[
 \|\widetilde h-\mathcal T_+ z_*\|_\infty
 \leq\|\widetilde h-\mathcal T_+\widetilde z\|_\infty
       +\|\mathcal T_+(\widetilde z-z_*)\|_\infty
 <3\delta+\eta.
\]
The final exact multiplication by $M$ therefore gives error less than
$M(3\delta+\eta)=\sqrt2\,M(3\ell+1)2^{-p}$, as claimed.
If one original operand has norm at most $\rho$, then
$\|\mathcal M_r(f,g)\|_\infty\leq\rho$.  Under the stated error
margin, the returned norm is less than $\rho+(1-\rho)=1$.

Every computed output before the final scaling is a disk grid array.
The layer's internal dyadic values and the integer product numerators
use the larger exact widths stated in their respective procedures.
The scale $M=2^{m_M}$ is a power of two, so it preserves the $p$-bit
grid. Compute $m_M=\sum_i\log_2t_i$ from the supplied axis lengths.
For each signed $w$-bit numerator returned by the opposite procedure,
append exactly $m_M$ zero bits at its least significant end. This is
the exact two's-complement encoding of the scaled numerator in
$w_M=w+m_M$ bits, including for negative values. The real word still
precedes the imaginary word and the coefficient order is unchanged.
Since $w\leq C_wp$ and $m_M\leq C_\ell p$, this known output width is
$O(p)$. Writing $\Gamma(w_M)$ takes $O(p)$ time.

There are $M=T/r$ pointwise polynomial products.  Their total cost
is $O(Tp\log(rp))$; the three transform costs are as stated, and
the last exact scaling is linear in $Tp$.  All data movement uses
the previously specified scans, swaps, and layout restoration.
\end{proof}
